\documentclass[11pt]{article}
\usepackage[a4paper, top=18mm, bottom=18mm, left=20mm, right=20mm]{geometry}
\usepackage[T2A]{fontenc}
\usepackage[utf8]{inputenc}
\usepackage[ukrainian]{babel}
\usepackage{graphicx}
\setlength{\parindent}{0pt}
\setlength{\parskip}{4pt}
\pagestyle{empty}
\begin{document}
%begin
{\large\textbf{Контрольна робота}}

\textbf{Варіант \No\,1}\hfill \textit{ПІБ:}\,\underline{\hspace{60mm}}
\vspace{4mm}

%beginitem
1. Що буде виведено за виконання фрагменту коду?
%Оригинал был утерян. Требует отладки!
%code
_a = 0

class A:
    _a = 1
    
    def __init__(self):
        self._a = 2
        _a = 3
        A._a = 4

obj = A()
obj._a = 5
try:
    print(_a, end=' ')
    print(obj._a, end=' ')
    print(A._a, end=' ')
    print(obj._a, end=' ')
    print(obj._A__a)
except:
    print('ERR')
%code

%enditem

%beginitem
2. Що буде виведено за виконання фрагменту коду?
9/9*9*9
%enditem

%beginitem
3. Що буде виведено за виконання фрагменту коду?

\begin{center}
	\adjustimage{max size={0.8\linewidth}{0.8\paperheight}}
	{../BinTree/trees_exp/110.png}
\end{center}
Указати послідовність міток вершин дерева, зображеного на рис., що утвориться в результаті обходу дерева в прямому порядку. ̳тки записувати через пробіл.\par Обхід дерева починається з кореня (на рис. це верхня вершина).
%answer:110 прямому
%enditem

%beginitem
4. %goodpath:Scripts/2018/Mod2018_1/01-good.txt
Відмітити все, що є однією правильною лексемою:
( 1 ) "1,3,2"

( 2 ) ||

( 3 ) 2int

( 4 ) [1]

%enditem

%beginitem
5. Що буде виведено за виконання фрагменту коду?
Записати ВИРАЗ, значенням якого є множина, що складається з кубівцілих чисел від 16 до 2006 (включно), що не діляться на жодне з чисел 2, 3.

%enditem

%beginitem
6. Що буде виведено за виконання фрагменту коду?

%code
a = 5
b = 1
c = 8

def h(a):
global c    a=3
    b*=2
    c=5
    return a+b+c

a, b, c = 2, 9, 4
try:
    print(h(b), a, b, c, sep=';')
except:
    print('Z', a, b, c, sep=';')
%code

%enditem

%beginitem
7. Що буде виведено за виконання фрагменту коду?

%code
f=[10,11,12,13]
f+=[2,3]
print(f)
%code


%enditem

%beginitem
8. Що буде виведено за виконання фрагменту коду?
Записати ВИРАЗ, значенням якого є множина, що складається з кубівцілих чисел від 16 до 2006 (включно), що не діляться на жодне з чисел 2, 3.
%enditem

%beginitem
9. Що буде виведено за виконання фрагменту коду?
f=('a','b','c','d','e',0,1,2,3); print(f[6])
%enditem

%beginitem
10. Що буде виведено за виконання фрагменту коду?

%code
#include <iostream>

int h(int a, int b){
    int c = 87;
    if (b != 1)
        return 5;
    if (a <= -4)
         c = 1;
    else 
        return 8;
    return c;
}

int main(){
    std::cout << h(-3, -3);
    return 0;
}
%code

%enditem

%beginitem
11. Що буде виведено за виконання фрагменту коду?
Рядки журналу через = містять ім'я користувача (ідентифікатор), час події,тип події, код події, наприклад
\begin{verbatim}
s="username = 2022-05-05 13:26 = ERROR = 404"
\end{verbatim}
у коді 
\begin{verbatim}
res = re.fullmatch('???', s) 
s = ''
code_str = ??? 
\end{verbatim}
чим треба замінити кожен з ???, щоб значенням code\_str був код події? 



%enditem

%beginitem
12. Що буде виведено за виконання фрагменту коду?
notTrue!=7.0 and not3.0<=7 and 3>=5
%enditem

%beginitem
13. Що буде виведено за виконання фрагменту коду?

%code
for f in range(-4,-4+4,-3):
    if f<8:
        continue
    print(f, end=' ')
    f=-4
    if f<8:
        break
else:
    print(f, end=' ')
print(f, end=' ')
%code

%enditem

%beginitem
14. Що буде виведено за виконання фрагменту коду?

%code
"True"!=7.0
%code

%enditem

%beginitem
15. Що буде виведено за виконання фрагменту коду?

%code
print(Trueis7.0is3.0is7)
%code

%enditem

%beginitem
16. Що буде виведено за виконання фрагменту коду?

%code
print(True is 7.0 <= 3.0 == 7)
%code

%enditem

%beginitem
17. Що буде виведено за виконання фрагменту коду?

%code
cout << (true != 7.0 <= 3.0 >= 7);
%code

%enditem

%beginitem
18. Що буде виведено за виконання фрагменту коду?

print('R', end=' ')
f=[10,11,12,13,14]
f+=[2,3]
print(f)


%enditem

%beginitem
19. %goodpath:Scripts/2019/Mod2019_1/Lex/03-good.txt
Відмітити все, що є коректним оператором С++:
( 1 ) print(3%50)

( 2 ) x<3 or y<7

( 3 ) (2+2j)%2

( 4 ) if ('x'><'y'): a=6

%enditem

%beginitem
20. Що буде виведено за виконання фрагменту коду?
a = 5
b = 1
c = 8

def h(a):
global c    a = 3
    b *= 2
    c = 5
    return a + b + c

a = 2
b = 9
c = 4

try:
    print(h(b), a, b, c, sep=';')
except:
    print('Z', a, b, c, sep=';')

%enditem

%beginitem
21. Що буде виведено за виконання фрагменту коду?

%code
#include <iostream>

int h(int a){
    int z = 87;
    if (a != 1) 
        return 5;
    if (a <= -4)
         z = 1;
    else
         return 8;
    return z;
}

int main(){
    std::cout << h(-3);
    return 0;
}
%code


%enditem

%beginitem
22. Що буде виведено за виконання фрагменту коду?

a,b,c=5,1,8
def h(a):
global c    a=3
    b*=2
    c=5
    return a+b+c

a,b,c=2,9,4
print(h(b),a,b,c)

%enditem

%beginitem
23. Що буде виведено за виконання фрагменту коду?
Записати ВИРАЗ, значенням якого є список, що складається з цілих чисел від 16 до 2006 (включно), причому спочатку за зростанням записано всі, що не діляться на 2, а потім решту за зростанням.

%enditem

%beginitem
24. Що буде виведено за виконання фрагменту коду?

%code
a = 5
b = 1
c = 8

def h(a):
global c    a=3
    b=2
    c=5
    return a+b+c

a, b, c = 2,9,4
try:
    print(h(b), a, b, c, sep=';')
except:
    print('Z', a, b, c, sep=';')
%code

%enditem

%beginitem
25. Що буде виведено за виконання фрагменту коду?

\begin{center}
	\adjustimage{max size={0.8\linewidth}{0.8\paperheight}}
	{../BinTree/trees_exp/110.png}
\end{center}
Указати послідовність міток вершин дерева, зображеного на рис., що утвориться в результаті обходу дерева в прямому порядку. Мітки записувати через пробіл.\par Алгоритм починає роботу з кореня (на рис. це верхня вершина).
%answer:110 прямому

%enditem

%beginitem
26. Що буде виведено за виконання фрагменту коду?
def h(a):
      z=87
      if a!=1: return 5
      if a<=-4: z=1
      else: return 8
      return z
   print(h(-3))
%enditem

%beginitem
27. Що буде виведено за виконання фрагменту коду?
9**0.5*False
%enditem

%beginitem
28. Що буде виведено за виконання фрагменту коду?

%code
print("True"!=7.0)
%code

%enditem

%beginitem
29. Що буде виведено за виконання фрагменту коду?
"True"!=7.0
%enditem

%beginitem
30. Що буде виведено за виконання фрагменту коду?
for f in range(-9,-9,-3):
    if f<8:
        break
        print(f,end=' ')
    if f<8:
        break
else:
    print(f, end=' ')
print(f, end=' ')

%enditem

%beginitem
31. Що буде виведено за виконання фрагменту коду?
Записати ВИРАЗ, значенням якого є список, що складається з кубівцілих чисел від 16 до 2006 (включно), що не діляться на жодне з чисел 2, 3, записаних в обернено\-му порядку.

%enditem

%beginitem
32. Що буде виведено за виконання фрагменту коду?

%code
try:
    t = 11
    while t>5:
        t += -3
    print(t, end=';')
except:
    print('Z')

%code

%enditem

%beginitem
33. Що буде виведено за виконання фрагменту коду?

print('R', end=' ')
t=5
while t<11:
    t+=1
print(t)


%enditem

%beginitem
34. Що буде виведено за виконання фрагменту коду?
Записати ВИРАЗ, значенням якого є словник, ключами якого є цілі числа від 16 до 2006 (включно), що не діляться на жодне з чисел 2, 3, а ключам відповідають їх квадратні корені 

%enditem

%beginitem
35. Що буде виведено за виконання фрагменту коду?

%code
print('R', end=' ')
for f in range(-4,-4,-1):
    if f<8:
        break
        print(f, end=' ')
    if f<8:
        break
    else:
        print(f, end=' ')
    print(f, end=' ')
%code

%enditem

%beginitem
36. Що буде виведено за виконання фрагменту коду?
try:
    for f in range(-9, -9, -3):
        if f<8:
            break
        print(f, end=';')
        if f<8:
            break
    else:
        print(f,end=';')
    print(f)
except:
    print('ERR')

%enditem

%beginitem
37. Що буде виведено за виконання фрагменту коду?

print('R', end=' ')
f=[0,1,2,3,4,5,6,7,8,9]
print(f[:-9:-3])


%enditem

%beginitem
38. Що буде виведено за виконання фрагменту коду?
Записати ВИРАЗ, значенням якого є список, що складається з цілих чисел від 16 до 2006 (включно), причому спочатку за зростанням записано всі, що не діляться на 2, а потім решту за зростанням.
%enditem

%beginitem
39. Що буде виведено за виконання фрагменту коду?
for f in range(-9,-9,-3):
    if f<8:
        break
        print(f,end=' ')
    if f<8:
        break
    else:
        print(f,end=' ')
    print(f)

%enditem

%beginitem
40. Що буде виведено за виконання фрагменту коду?
try:
    res = "True"!=7.0
    if isinstance(res, bool):
        print(f'bool;{res}')
    elif isinstance(res, int):
        print(f'int;{res}')
    elif isinstance(res, float):
        print(f'float;{res:.2f}')
    else:
        print('???')
except:
    print('ERR')

%enditem

%beginitem
41. Що буде виведено за виконання фрагменту коду?

%code
print(9/9*9*9)
%code

%enditem

%beginitem
42. Що буде виведено за виконання фрагменту коду?

%code
#include <iostream>
using namespace std;

int a = 5, b = 1, c = 8;

int h(int &a){
int c;    a = 3;
    b -= 2;
    c = 5;
    return a + b + c;
}

int main(){
    inta = 2;
    intb = 9;
    intc = 4;
    cout << h(b) << ':';
    cout << a << ':' << b << ':' << c;
    return 0; 
}
%code

%enditem

%beginitem
43. Що буде виведено за виконання фрагменту коду?

%code
a,b,c=5,1,8
def h(a):
    a=3
    b*=2
    c=5
    return a+b+c

a,b,c=2,9,4
print(h(b),a,b,c)
%code

%enditem

%beginitem
44. Що буде виведено за виконання фрагменту коду?
%code

def f(n):
    print("f", end="");
    return n

print(f(-3) and f(-8))
%code

%enditem

%beginitem
45. Що буде виведено за виконання фрагменту коду?
f=[10,11,12,13,14]; f+=[2,3]; print(f)
%enditem

%beginitem
46. Що буде виведено за виконання фрагменту коду?
Написати програму, що запитує у користувача значення дійсних параметрів $x$ та $y$, обчислює для них значення виразу 
$$\frac{\log_{2} x + 50}{(\arctg x + 0.7) (y + 63)}$$ 
та повiдомляє введенi данi та результат обчислення.


%enditem

%beginitem
47. Що буде виведено за виконання фрагменту коду?

print('R', end=' ')
for f in range(-9,-9,-3):
    if f<8:
        break
        print(f, end=' ')
    if f<8:
        break
    else:
        print(f, end=' ')
    print(f)

%enditem

%beginitem
48. Що буде виведено за виконання фрагменту коду?

%code
f = ('a','b','c','d','e',0,1,2,3)
print(f[6])
%code


%enditem

%beginitem
49. Що буде виведено за виконання фрагменту коду?
В умовах попередньої задачі було виконано p->prev=p->next->next;

Чому дорівнює значення p->prev->prev->prev->prev->n ?
%enditem

%beginitem
50. Що буде виведено за виконання фрагменту коду?
Скільки 2017-елементних підмножин множини натуральних від 1 до 20172017 містять принаймні одне число, що більше 172005
%enditem

%beginitem
51. Що буде виведено за виконання фрагменту коду?
try:
    t = 11
    while t>5:
        t += -3
    print(t, end=';')
except:
    print('Z')

%enditem

%beginitem
52. Що буде виведено за виконання фрагменту коду?
Записати ВИРАЗ, значенням якого є множина, що складається з цілих чисел від 16 до 2006 (включно), що не діляться на 2.
%enditem

%beginitem
53. Що буде виведено за виконання фрагменту коду?
Задано тип вузла списку
struct Node \{int n; Node *prev, *next;\};
У списку зберігаються послідовні цілі числа від -45 до 40. Вказівник p встановлено на вузол, в якому зберігається число -3.
Чому дорівнює значення p->prev->prev->prev->prev->prev->n ?
%enditem

%beginitem
54. %goodpath:Scripts/2018/Mod2018_1/03-good.txt
Відмітити все, що є коректним оператором С++:
( 1 ) if ('x'<'y') a=6;

( 2 ) x<y<z

( 3 ) if ('x'<>'y') a=6;

( 4 ) cin<<x;

%enditem

%beginitem
55. Що буде виведено за виконання фрагменту коду?
try:
    res = 9/9*9*9
    if isinstance(res, bool):
        print(f'bool;{res}')
    elif isinstance(res, int):
        print(f'int;{res}')
    elif isinstance(res, float):
        print(f'float;{res:.2f}')
    else:
        print('???')
except:
    print('Z')

%enditem

%beginitem
56. Що буде виведено за виконання фрагменту коду?
Знайти кількість відношень на n-елементній множині, які нерефлексивні та несиметричні
%enditem

%beginitem
57. Що буде виведено за виконання фрагменту коду?

print('R', end=' ')
t=11
while t>5:
    t+=-3
print(t)


%enditem

%beginitem
58. Що буде виведено за виконання фрагменту коду?
$a$ є цілочисловим масивом типу $numpy.ndarray$ розмірів $2{\times}2$. Сума елементів масиву $a$ дорівнює 26. Чому дорівнює сума елементів масиву $a$ після виконання 
\begin{verbatim}
a += 26
\end{verbatim}


Якщо код некоректний, то в якості відповіді зазначити ERROR.



%enditem

%beginitem
59. Що буде виведено за виконання фрагменту коду?
f=[10,11,12,13,14]; del f[:-4]; print(f)
%enditem

%beginitem
60. Що буде виведено за виконання фрагменту коду?

%code
print('R', end=' ')
f=[10,11,12,13,14]
del f[:-4]
print(f)
%code


%enditem

%beginitem
61. Що буде виведено за виконання фрагменту коду?
Записати ВИРАЗ, значенням якого є список, що складається з цілих чисел від 16 до 2006 (включно), причому спочатку за зростанням записано всі, що не діляться на 2, а потім решту за зростанням.
%enditem

%beginitem
62. Що буде виведено за виконання фрагменту коду?
%code
for f in range(-9,-9,-3):
    if f<8:
        break
        print(f,end=' ')
    if f<8:
        break
else:
    print(f, end=' ')
print(f, end=' ')
%code

%enditem

%beginitem
63. Що буде виведено за виконання фрагменту коду?

def f(n):
    print("f", end="");
    return n

print(f(-3) and f(-8))

%enditem

%beginitem
64. Що буде виведено за виконання фрагменту коду?

print('R', end=' ')
t=5
while t<11:
    t+=1
print(t)


%enditem

%beginitem
65. Що буде виведено за виконання фрагменту коду?
t=11
   while t>5: t+=-3
   print(t)
%enditem

%beginitem
66. Що буде виведено за виконання фрагменту коду?

%code
#include <iostream>
using namespace std;

int h(int a, int b){
    int c = 87;
    if (b != 1)
        return 5;
    if (a <= -4)
         c = 1;
    else 
        return 8;
    return c;
}

int main(){
    cout << h(-3, -3);
    return 0;
}
%code

%enditem

%beginitem
67. Що буде виведено за виконання фрагменту коду?
a,b,c=5,1,8
   def h(a):
     global c     a=3
     b=2
     c=5
     return a+b+c
   a,b,c=2,9,4
   print(h(b),a,b,c)
%enditem

%beginitem
68. Що буде виведено за виконання фрагменту коду?

%code
f=[10,11,12,13]
del f[:-4]
print(f)
%code


%enditem

%beginitem
69. Що буде виведено за виконання фрагменту коду?
def h(a,b):
      c=87
      if b!=1: return 5
      if a<=-4: c=1
      else: return 8
      return c
   print(h(-3,-3))
%enditem

%beginitem
70. %goodpath:Scripts/2019/Mod2019_1/Lex/01-good.txt
Відмітити все, що є однією правильною лексемою:
( 1 ) Switch

( 2 ) **

( 3 ) double

( 4 ) -56

%enditem

%beginitem
71. Що буде виведено за виконання фрагменту коду?

\begin{center}
	\adjustimage{max size={0.8\linewidth}{0.8\paperheight}}
	{../BinTree/trees_exp/110.png}
\end{center}
Указати послідовність міток вершин дерева, зображеного на рис., що утвориться в результаті обходу дерева в прямому порядку. Мітки записувати через пробіл.\par Обхід дерева починається з кореня (на рис. це верхня вершина).
%answer:110 прямому

%enditem

%beginitem
72. Що буде виведено за виконання фрагменту коду?
Змінні \kw{next} та \kw{prev} вузла списку позначають наступний та попередній за ним вузол відповідно. Починаючи з голови, у список записано послідовні цілі числа від~$-45$ до~$40$. Нехай змінна \kw{p} позначає вузол, що зберігає значення~$-3$.
Яке число зберігається у вузлі \kw{p.prev.prev.prev.prev.prev} ?

%enditem

%beginitem
73. Що буде виведено за виконання фрагменту коду?

print('R', end=' ')
f=[0,1,2,3,4,5,6,7,8,9]
print(f[:-9:-3])


%enditem

%beginitem
74. Що буде виведено за виконання фрагменту коду?

print('R', end=' ')
f=[10,11,12,13,14]
del f[:-4]
print(f)


%enditem

%beginitem
75. Що буде виведено за виконання фрагменту коду?
Рядки журналу через двокрапку містять ім'я користувача (ідентифікатор), час події,тип події, код події, наприклад
\begin{verbatim}
s="username : 2022-05-05 13:26 : ERROR : 404"
\end{verbatim}
у коді 
\begin{verbatim}
res = re.fullmatch('???', s) 
s = ''
code_str = ??? 
\end{verbatim}
чим треба замінити кожен з ???, щоб значенням code\_str був код події? 



%enditem

%beginitem
76. Що буде виведено за виконання фрагменту коду?
try:
    t = 5
    while t<11:
        t += 1
    print(t, end=';')
except:
    print('Z')

%enditem

%beginitem
77. Що буде виведено за виконання фрагменту коду?
Записати ВИРАЗ, значенням якого є множина, що складається з цілих чисел від 16 до 2006 (включно), що не діляться на 2.
%enditem

%beginitem
78. Що буде виведено за виконання фрагменту коду?

a,b,c=6,7,9
def h(a,b,c):
    print(a,b,c,end="")

h(5,c=1,b=3)
print(a,b,c)


%enditem

%beginitem
79. Що буде виведено за виконання фрагменту коду?
Написати програму, що запитує у користувача значення дійсних параметрів $x$ та $y$, обчислює для них значення виразу та повiдомляє введенi данi та результат обчислення.
$$\frac{\log_{2} x + 50}{(\arctg x + 0.7) (y + 63)}$$ 


%enditem

%beginitem
80. Що буде виведено за виконання фрагменту коду?
"True"!=7.0
%enditem

%beginitem
81. Що буде виведено за виконання фрагменту коду?
try:
    res = 9**0.5*False
    if isinstance(res, bool):
        print(f'bool;{res}')
    elif isinstance(res, int):
        print(f'int;{res}')
    elif isinstance(res, float):
        print(f'float;{res:.2f}')
    else:
        print('???')
except:
    print('ERR')

%enditem

%beginitem
82. Що буде виведено за виконання фрагменту коду?
%code
if (5 != 5)
    cout << "z";
else
    cout << "b";
%code


%enditem

%beginitem
83. Що буде виведено за виконання фрагменту коду?
def f(n):
    print('f', end='')
    return n

try:
    print(f(-3) and f(-8))
except:
    print('ERR')

%enditem

%beginitem
84. Що буде виведено за виконання фрагменту коду?

%code
cout << 10 % 10 % 10 % 10;
%code

%enditem

%beginitem
85. Що буде виведено за виконання фрагменту коду?
try:
    print(5, end=';')
    print(1j!=5j, end=';')
    print(8, end=';')
except TypeError:
    print(2, end=';')
except ValueError:
    print(9, end=';')
except:
    print('ERR', end=';')
else:
    print(4, end=';')
finally:
    print(6)

%enditem

%beginitem
86. Що буде виведено за виконання фрагменту коду?

print('R', end=' ')
t=11
while t>5:
    t+=-3
print(t)


%enditem

%beginitem
87. Що буде виведено за виконання фрагменту коду?

%code
a,b,c=5,1,8
def h(a):
global c    a=3
    b=2
    c=5
    return a+b+c

a,b,c=2,9,4
print(h(b),a,b,c)
%code

%enditem

%beginitem
88. Що буде виведено за виконання фрагменту коду?

%code
cout << (!true!=7.0 and !3.0<= 7 and 3>=5);
%code

%enditem

%beginitem
89. Що буде виведено за виконання фрагменту коду?

%code
#include <iostream>
using namespace std;

bool f(int n){
    cout<<"f";
    return n;
}

int main(){
    cout<<(f(-3) and f(-8));
    return 0;
}
%code

%enditem

%beginitem
90. Що буде виведено за виконання фрагменту коду?

%code
cout << (9 % 9 % 9 % 9);
%code

%enditem

%beginitem
91. Що буде виведено за виконання фрагменту коду?
f=[10,11,12,13,14]; del f[:-4]; print(f)
%enditem

%beginitem
92. Що буде виведено за виконання фрагменту коду?

%code
print(9**0.5*False)
%code

%enditem

%beginitem
93. Що буде виведено за виконання фрагменту коду?
Записати ВИРАЗ, значенням якого є список, що складається з кубівцілих чисел від 16 до 2006 (включно), що не діляться на жодне з чисел 2, 3, записаних в оберненому порядку.
%enditem

%beginitem
94. Що буде виведено за виконання фрагменту коду?

a,b,c=5,1,8
def h(a):
global c    a=3
    b=2
    c=5
    return a+b+c

a,b,c=2,9,4
print(h(b),a,b,c)

%enditem

%beginitem
95. Що буде виведено за виконання фрагменту коду?

\begin{center}
	\adjustimage{max size={0.8\linewidth}{0.8\paperheight}}
	{../15BinTree/trees_exp/110.png}
\end{center}
Указати послідовність міток вершин дерева, зображеного на рис., що утвориться в результаті обходу дерева в прямому порядку. Мітки записувати через пробіл.\par Алгоритм починає роботу з кореня (на рис. це верхня вершина).
%answer:110 прямому

%enditem

%beginitem
96. Що буде виведено за виконання фрагменту коду?
Рядки журналу через = містять ім'я користувача (ідентифікатор), час події,тип події, код події, наприклад
\begin{verbatim}
s="username = 2022-05-05 13:26 = ERROR = 404"
\end{verbatim}
у коді 
\begin{verbatim}
res = re.fullmatch('???', s) 
s = ''
code_str = ??? 
\end{verbatim}
чим треба замінити кожен з ???, щоб значенням code\_str був код події? 



%enditem

%beginitem
97. Що буде виведено за виконання фрагменту коду?
Задано тип вузла списку struct Node {int n; Node *prev, *next;};
У списку зберігаються послідовні цілі числа від -45 до 40. Вказівник p встановлено на вузол, в якому зберігається число -3.
Чому дорівнює значення p->prev->prev->prev->prev->prev->n ?
%enditem

%beginitem
98. Що буде виведено за виконання фрагменту коду?

%code
print('R', end=' ')
f = ('a','b','c','d','e',0,1,2,3)
print(f[6])
%code


%enditem

%beginitem
99. Що буде виведено за виконання фрагменту коду?
%code
#include <iostream>
using namespace std;

int f(int &x, int &y){
    x = 7;
    y+= 2;
    return x;
}

int main(){
    int a = 5, b = 4;
    a = f(a, a);
    cout << a << ":" << b <<':';
    {
        int a = 1, b = 9;
        cout << ((a<5) && ((b+=1) < 5)) << ':';
        cout << a << ":" << b << ':';
    }
    {
        inta = 6;
        cout << a << ':';
    }
    cout << a << ":" << b << endl;
    return 0;
}
%code


%enditem

%beginitem
100. Що буде виведено за виконання фрагменту коду?

%code
try:
    t = 5
    while t<11:
        t += 1
    print(t, end=';')
except:
    print('Z')

%code

%enditem

%end
\newpage
\end{document}

